\section{Our Contributions}
\label{sec:our-contributions}

\paragraph{Positioning.}
Rosetta, Grafite, and Memento all carry a worst-case FPR guarantee ---
the Goswami bound up to a constant or additive term --- and differ
mainly in cost: Rosetta pays $O(\log_2 \mathcal{L})$ Bloom probes per
query, Grafite a single predecessor probe within an additive $2n$ bits,
and Memento adds dynamicity at static-axis parity with Grafite.
SuRF and SNARF trade that guarantee for shape: SuRF is not
parameterized by $\varepsilon$ at all and inherits its FPR from the
trie skeleton (empirically up to $25\%$ on correlated workloads);
SNARF goes below the Goswami bound when its CDF model fits the data
and loses the guarantee otherwise.
None of the five exploits the cluster-and-gap structure common in real
LSM workloads.
We keep the worst-case guarantee of the Grafite line and add a
distribution-aware mechanism that beats the bound on realistic inputs.

\paragraph{Contributions of this thesis.}
We start from Goswami's two-layer construction --- locality-preserving
hash followed by an ERE backend --- and extend it in multiple areas,
producing a filter that matches the Goswami
lower bound on adversarial inputs and beats it on realistic ones.
\begin{enumerate}
    \item \textbf{Single-vector ERE backend (Chapter~\ref{chap:ere}).}
    We adopt the single-vector $D$ layout (also used by
    Grafite~\citep{grafite2024}) in place of Goswami's two metadata
    vectors $D_1, D_2$, and back it with a Poisson bucket-occupancy
    model that predicts the resulting space saving --- $24\%$ relative
    to the vanilla Goswami backend, confirmed empirically
    (Section~\ref{sec:motivation} sketches the empirical reasoning
    that justifies this collapse).

    \item \textbf{Adaptive bucket search (Chapter~\ref{chap:ere}).}
    We implement Goswami's Weak Prefix Search (WPS) index in full and
    show empirically that adaptive binary/linear search on packed
    suffixes outperforms it at LSM bucket sizes, removing $59$~bits of
    auxiliary metadata per key and reducing per-query latency by
    $13$--$30\times$ relative to WPS.

    \item \textbf{Truncation hash (Chapter~\ref{chap:hash}).}
    For sparse inputs we replace Goswami's locality-preserving
    rotation with the truncation map $h(x) = \lfloor x/2^t \rfloor$,
    which saves $\log_2 \mathcal{L}$ bits per key.

    \item \textbf{Distribution-aware segmentation: Seg-ARE (Chapter~\ref{chap:hybrid}).}
    A single-pass density-based segmentation partitions the input
    into dense clusters --- handed to exact-mode ERE sub-filters
    (FPR $= 0$) --- and sparse tails routed to the truncation hash.
    On SOSD dataset this drives the FPR below $10^{-8}$ at
    $\sim\!11$ bits per key on Facebook User IDs and $\sim\!6$ bits per key on
    Wiki timestamps, matching or beating the strongest baselines
    (Chapter~\ref{chap:evaluation}); Chapter~\ref{chap:conclusion}
    discusses when Seg-ARE wins, ties, and loses, and the
    Elias--Fano interpretation behind it.

    \item \textbf{Improvements to the Grafite reference implementation.}
    Beyond designing our own filter, we deeply analysed the closest
    prior art in the Goswami line and patched it with two
    fixes applied in every measurement of
    Chapter~\ref{chap:evaluation}: a balanced-descent replacement for
    the SUX \texttt{select\_zero} pathology on clustered inputs
    (three-orders-of-magnitude latency regression in the unpatched
    code, Section~\ref{subsec:grafite-improvements}); and a lossless
    Elias--Fano fallback that activates when the upstream constructor
    would otherwise throw on configurations where lossless storage is
    cheaper than the requested approximate filter
    (Section~\ref{subsec:grafite-fallback}). Grafite is therefore
    reported in its best configuration throughout this work.
\end{enumerate}

\paragraph{Roadmap.}
\begin{itemize}
    \item \textbf{Chapter~\ref{chap:ere}} develops the single-vector
        ERE and the adaptive bucket-search backend
        (contributions~1--2).
    \item \textbf{Chapter~\ref{chap:distributions}} characterises the
        key distributions we target, preparing the ground for
        contributions~3--4.
    \item \textbf{Chapter~\ref{chap:hash}} formalises and compares the
        SODA and Truncation hash designs (contribution~3).
    \item \textbf{Chapter~\ref{chap:hybrid}} introduces the Seg-ARE
        segmentation (contribution~4).
    \item \textbf{Chapter~\ref{chap:evaluation}} reports the empirical
        comparison against SuRF, SNARF, Rosetta, and Grafite on real world SOSD
        and synthetic workloads.
    \item \textbf{Chapter~\ref{chap:conclusion}} concludes.
\end{itemize}
